% =====================================================================
%  tikzlib · ciencias/lewis — puntos de Lewis, enlaces y esquemas de enlace
%  Autor: lote videos 6-7 (química 9º). Recurso ORIGINAL (sin dependencias
%  externas) → apto para exportar a SVG (pdftocairo) y animar en Manim.
%  Requiere: \usepackage{tikz} + \usetikzlibrary{calc}
%
%  PRIMITIVAS
%    \lewisdot{x}{y}                 un electrón (punto lleno) en (x,y)
%    \lewissym{S}{n}                 símbolo de Lewis: S con n e- de valencia
%                                    (1..8), reparte 1 por lado y luego aparea
%
%  MOLÉCULAS (dentro de una tikzpicture con nodos ya colocados)
%    \lonepair{ang}{nodo}           par libre: 2 puntos a distancia, en 'ang°'
%    \lwsingle{A}{B}                enlace simple  A—B
%    \lwdouble{A}{B}                enlace doble   A=B   (bonos horizontales)
%    \lwtriple{A}{B}                enlace triple  A≡B   (bonos horizontales)
%    \shpair{A}{B}                  par compartido: 2 puntos en el punto medio
%
%  ESQUEMAS DE ENLACE (auto-contenidos: dibújalos en su tikzpicture)
%    \marelectrones                 "mar de electrones" (metálico): rejilla de
%                                   cationes + e- libres dispersos
%    \redionica                     red iónica NaCl (cationes/aniones alternos)
%
%  Preview: se exportan versiones SVG en assets/svg-src/quimica/
% =====================================================================

% --- un electrón ------------------------------------------------------
\newcommand{\lewisdot}[2]{\fill[ink] (#1,#2) circle (1.7pt);}

% distancia del punto al centro del símbolo y media-separación del par
\def\lwR{0.44}
\def\lwP{0.15}

% --- símbolo de Lewis: S con n electrones de valencia (1..8) ----------
\newcommand{\lewissym}[2]{%
\begin{tikzpicture}[baseline=(sym.base)]
  \node[font=\sffamily\large\bfseries, inner sep=1pt, text=ink] (sym) at (0,0) {#1};
  \pgfmathtruncatemacro{\lwn}{#2}%
  % derecha (electrones 1 y 5)
  \ifnum\lwn>4 \lewisdot{\lwR}{\lwP}\lewisdot{\lwR}{-\lwP}%
  \else\ifnum\lwn>0 \lewisdot{\lwR}{0}\fi\fi
  % arriba (2 y 6)
  \ifnum\lwn>5 \lewisdot{\lwP}{\lwR}\lewisdot{-\lwP}{\lwR}%
  \else\ifnum\lwn>1 \lewisdot{0}{\lwR}\fi\fi
  % izquierda (3 y 7)
  \ifnum\lwn>6 \lewisdot{-\lwR}{\lwP}\lewisdot{-\lwR}{-\lwP}%
  \else\ifnum\lwn>2 \lewisdot{-\lwR}{0}\fi\fi
  % abajo (4 y 8)
  \ifnum\lwn>7 \lewisdot{\lwP}{-\lwR}\lewisdot{-\lwP}{-\lwR}%
  \else\ifnum\lwn>3 \lewisdot{0}{-\lwR}\fi\fi
\end{tikzpicture}}

% --- par libre sobre un nodo, a distancia, en dirección 'ang' grados --
% Dos puntos que flanquean la línea radial (perpendicular al radio).
\newcommand{\lonepair}[2]{%
  \coordinate (lpc) at ($(#2)+(#1:0.46)$);
  \fill[ink] ($(lpc)+(#1+90:0.10)$) circle (1.6pt);
  \fill[ink] ($(lpc)+(#1-90:0.10)$) circle (1.6pt);}

% --- enlaces (bonos horizontales para dobles/triples) -----------------
% transform canvas conserva el recorte en el borde del nodo (el enlace queda
% en el HUECO entre las letras) y separa las líneas paralelas del doble/triple.
\newcommand{\lwsingle}[2]{\draw[line width=1.1pt, ink] (#1)--(#2);}
\newcommand{\lwdouble}[2]{%
  \draw[line width=1.1pt, ink, transform canvas={yshift=3pt}] (#1)--(#2);
  \draw[line width=1.1pt, ink, transform canvas={yshift=-3pt}] (#1)--(#2);}
\newcommand{\lwtriple}[2]{%
  \draw[line width=1.1pt, ink, transform canvas={yshift=5pt}] (#1)--(#2);
  \draw[line width=1.1pt, ink] (#1)--(#2);
  \draw[line width=1.1pt, ink, transform canvas={yshift=-5pt}] (#1)--(#2);}

% --- par de electrones compartido (2 puntos en el punto medio A-B) ----
\newcommand{\shpair}[2]{%
  \coordinate (shm) at ($(#1)!0.5!(#2)$);
  \fill[ink] ($(shm)+(0,0.08)$) circle (1.6pt);
  \fill[ink] ($(shm)+(0,-0.08)$) circle (1.6pt);}

% =====================================================================
%  ESQUEMA · MAR DE ELECTRONES (enlace metálico)
%  Rejilla de cationes (+) fijos + electrones deslocalizados dispersos.
% =====================================================================
% Rellenos PLANOS (sin gradientes) → aptos para exportar a SVG/SVGMobject.
\newcommand{\marelectrones}{%
\begin{tikzpicture}[scale=1]
  % cationes en rejilla 4x3
  \foreach \x in {0,1,2,3}{
    \foreach \y in {0,1,2}{
      \fill[subjectcolor!70] (\x*1.15,\y*1.15) circle (0.30);
      \draw[subjectcolor!85, line width=0.6pt] (\x*1.15,\y*1.15) circle (0.30);
      \node[font=\sffamily\bfseries\small, text=white] at (\x*1.15,\y*1.15) {+};
    }}
  % electrones deslocalizados (entre los cationes)
  \foreach \p in {(0.55,0.55),(1.75,0.4),(2.9,0.7),(0.4,1.75),(1.65,1.6),
                  (2.75,1.8),(3.4,1.2),(1.1,1.15),(2.25,1.15),(0.55,-0.05),
                  (2.3,0.05),(3.5,2.3)}{
    \fill[cnatlima] \p circle (0.085);}
\end{tikzpicture}}

% =====================================================================
%  ESQUEMA · RED IÓNICA (tipo NaCl)  cationes (+) y aniones (-) alternos
% =====================================================================
\newcommand{\redionica}{%
\begin{tikzpicture}[scale=1]
  \foreach \x in {0,1,2,3}{
    \foreach \y in {0,1,2,3}{
      \pgfmathtruncatemacro{\sum}{\x+\y}%
      \ifodd\sum
        \fill[cnatlima!80] (\x*1.05,\y*1.05) circle (0.34)
          node[font=\sffamily\bfseries\scriptsize, text=ink]{Cl$^-$};
      \else
        \fill[subjectcolor!75] (\x*1.05,\y*1.05) circle (0.24)
          node[font=\sffamily\bfseries\tiny, text=white]{Na$^+$};
      \fi
    }}
  % líneas de red tenues
  \draw[subjectcolor!25, line width=0.5pt] (0,0) grid[step=1.05] (3*1.05,3*1.05);
\end{tikzpicture}}
